\chapter{Corporate Finance Foundations, Aging Analysis \& Outlier Detection}

% ----------------------------------------------------------------------
% PART 1: EXECUTIVE ROADMAP & LEARNING OBJECTIVES
% ----------------------------------------------------------------------
\section{Executive Roadmap \& Learning Objectives}

Corporate financial management governs the allocation of scarce capital over time, anchored in the Time Value of Money (TVM) and rigorous capital budgeting hurdle rates. In forensic accounting and internal audit investigations, financial risk is frequently concentrated within working capital cycles---most notably within Accounts Receivable aging schedules and procurement disbursement streams. To detect anomalous transaction clusters, forensic data scientists must combine financial accounting rules with robust statistical outlier detection algorithms.

\begin{tcolorbox}[colback=primarylight, colframe=primarydark, title={\textbf{Core Learning Competencies}}]
Upon mastering this chapter, students will be equipped to:
\begin{itemize}
    \item \textbf{Evaluate Capital Budgeting Invariants:} Calculate Present Value ($PV$), Net Present Value ($NPV$), and Internal Rate of Return ($IRR$) to audit project cash flows and unmask inflated capitalization projections.
    \item \textbf{Construct Accounts Receivable Aging Matrices:} Stratify customer receivables across temporal aging buckets ($0$--$30$, $31$--$60$, $61$--$90$, $>90$ days) and compute expected credit loss allowances.
    \item \textbf{Apply Pareto 80/20 Stratification:} Isolate high-impact transactions through cumulative ABC spending stratification while auditing low-value anomaly tails.
    \item \textbf{Compare Outlier Detection Methodologies:} Contrast parametric $Z$-scores, non-parametric Tukey Interquartile Ranges ($IQR$), and robust Modified $Z$-scores based on Median Absolute Deviation ($MAD$).
    \item \textbf{Deploy Automated Anomaly Detectors:} Program Python engines to ingest transaction journals, compute aging latencies, and flag multivariate statistical outliers.
\end{itemize}
\end{tcolorbox}

\begin{table}[htbp]
\centering
\small
\begin{tabularx}{\textwidth}{llX}
\toprule
\textbf{Construct / Metric} & \textbf{Symbol} & \textbf{Mathematical Definition / Formula} \\
\midrule
Present Value & $PV$ & $\frac{FV}{(1 + r/m)^{m \cdot t}} \quad \text{or} \quad FV \cdot e^{-rt}$ (continuous) \\
Net Present Value & $NPV$ & $\sum_{t=0}^n \frac{CF_t}{(1 + k)^t} = -CF_0 + \sum_{t=1}^n \frac{CF_t}{(1 + k)^t}$ \\
Internal Rate of Return & $IRR$ & Discount rate setting $NPV \equiv 0$: $\sum_{t=0}^n \frac{CF_t}{(1 + IRR)^t} = 0$ \\
Standard $Z$-Score & $Z_i$ & $\frac{x_i - \bar{x}}{s} \quad \text{(sensitive to extreme values)}$ \\
Tukey IQR Boundaries & $[L_B, U_B]$ & $[Q_1 - 1.5 \cdot IQR, \; Q_3 + 1.5 \cdot IQR] \quad \text{where } IQR = Q_3 - Q_1$ \\
Median Absolute Deviation & $MAD$ & $\text{median}\left(|x_i - \text{median}(x)|\right)$ \\
Modified $Z$-Score & $M_i$ & $\frac{0.6745 \cdot (x_i - \text{median}(x))}{MAD} \quad (|M_i| > 3.5 \implies \text{Outlier})$ \\
\bottomrule
\end{tabularx}
\caption{Mathematical Notation Guide for Valuation, Aging, and Statistical Outliers.}
\label{tab:w6_notation}
\end{table}

% ----------------------------------------------------------------------
% PART 2: MATHEMATICAL & COMPUTATIONAL FOUNDATIONS
% ----------------------------------------------------------------------
\section{Mathematical \& Computational Foundations}

\subsection{The Time Value of Money and Discounted Cash Flows}
The fundamental postulate of corporate finance asserts that a monetary unit received today commands greater economic value than an identical unit promised in the future. For discrete compounding at nominal rate $r$ compounded $m$ times annually over $t$ years:
\begin{equation}
FV_t = PV_0 \left(1 + \frac{r}{m}\right)^{m \cdot t} \iff PV_0 = FV_t \left(1 + \frac{r}{m}\right)^{-m \cdot t}
\label{eq:tvm_discrete}
\end{equation}
In the continuous compounding limit ($m \to \infty$):
\begin{equation}
\lim_{m \to \infty} \left(1 + \frac{r}{m}\right)^{m \cdot t} = e^{rt} \implies PV_0 = FV_t \cdot e^{-rt}
\label{eq:tvm_continuous}
\end{equation}

\subsection{Capital Budgeting Hurdle Rates: NPV vs.~IRR}
Enterprises evaluate capital allocation through Net Present Value ($NPV$), discounting project net cash flows at the enterprise hurdle rate $k$:
\begin{equation}
NPV(k) = -CF_0 + \sum_{t=1}^n \frac{CF_t}{(1 + k)^t}
\label{eq:npv_formula}
\end{equation}
The Internal Rate of Return ($IRR$) is the unique root of the polynomial equation $NPV(IRR) = 0$. In capital budgeting forensics, project sponsors often manipulate forward cash flow estimates to push projected $IRR$ above internal hurdle rates, justifying fraudulent or unviable capital expenditure allocations.

\subsection{Statistical Outlier Detection Frameworks}
In enterprise ledgers, transactions can be evaluated through three distinct statistical paradigms:
\begin{enumerate}
    \item \textbf{Parametric Gaussian $Z$-Score:}
    \begin{equation}
    Z_i = \frac{x_i - \bar{x}}{s}, \quad \text{where } \bar{x} = \frac{1}{n} \sum_{i=1}^n x_i, \quad s = \sqrt{\frac{1}{n-1} \sum_{i=1}^n (x_i - \bar{x})^2}
    \label{eq:z_score}
    \end{equation}
    Under Gaussian assumptions, observations with $|Z_i| > 3.0$ occur with probability $p \approx 0.0027$.
    \item \textbf{Non-Parametric Tukey Boxplot (IQR):}
    \begin{equation}
    \text{Inner Fences} = [Q_1 - 1.5 \cdot IQR, \; Q_3 + 1.5 \cdot IQR]
    \label{eq:tukey_iqr}
    \end{equation}
    \begin{equation}
    \text{Outer Fences} = [Q_1 - 3.0 \cdot IQR, \; Q_3 + 3.0 \cdot IQR]
    \end{equation}
    Values beyond the outer fences represent severe outliers requiring substantive voucher testing.
    \item \textbf{Boris Iglewicz and David Hoaglin Modified $Z$-Score:}
    To prevent extreme fraudulent transactions from masking other outliers by inflating $\bar{x}$ and $s$, the Modified $Z$-score replaces sample moments with order statistics:
    \begin{equation}
    M_i = \frac{0.6745 \cdot (x_i - \tilde{x})}{MAD}, \quad \tilde{x} = \text{median}(x), \quad MAD = \text{median}(|x_i - \tilde{x}|)
    \label{eq:modified_z}
    \end{equation}
    The factor $0.6745$ ensures asymptotic equivalence to the standard deviation under a normal distribution because $\Phi(0.6745) = 0.75 \implies \text{P}(|X - \mu| \le 0.6745 \sigma) = 0.50$.
\end{enumerate}

% ----------------------------------------------------------------------
% PART 3: VISUAL & GEOMETRIC INTUITION
% ----------------------------------------------------------------------
\section{Visual \& Geometric Intuition: The Accounts Receivable Aging Risk Ladder}

Figure~\ref{fig:ar_aging_ladder} visualizes the temporal degradation of accounts receivable and the escalating loss provisioning rates required as invoices age past due.

\begin{figure}[htbp]
\centering
\begin{tikzpicture}[
    font=\sffamily\footnotesize,
    node distance=1.1cm and 1.4cm,
    step1/.style={rectangle, draw=tealborder, fill=tealbg, thick, rounded corners=3pt, minimum width=2.8cm, minimum height=1.1cm, align=center},
    step2/.style={rectangle, draw=primarydark, fill=primarylight, thick, rounded corners=3pt, minimum width=2.8cm, minimum height=1.1cm, align=center},
    step3/.style={rectangle, draw=amberborder, fill=amberbg, thick, rounded corners=3pt, minimum width=2.8cm, minimum height=1.1cm, align=center},
    step4/.style={rectangle, draw=coralborder, fill=coralbg, thick, rounded corners=3pt, minimum width=2.8cm, minimum height=1.1cm, align=center},
    arr/.style={->, >=stealth, thick, color=primarydark}
]

% Steps
\node[step1] (b1) {\textbf{Current (0--30 Days)}\\Default Rate: $1\%$--$2\%$\\Normal Trade Credit};
\node[step2, right=of b1] (b2) {\textbf{Past Due (31--60)}\\Default Rate: $5\%$--$10\%$\\Payment Delay Flags};
\node[step3, right=of b2] (b3) {\textbf{Past Due (61--90)}\\Default Rate: $20\%$--$35\%$\\Credit Hold Trigger};
\node[step4, right=of b3] (b4) {\textbf{Severe ($>90$ Days)}\\Default Rate: $50\%$--$100\%$\\Impairment Write-Off};

% Connectors
\draw[arr] (b1) -- (b2);
\draw[arr] (b2) -- (b3);
\draw[arr] (b3) -- (b4);

% Forensic Callout
\node[draw=red, dashed, fill=red!5, rounded corners=3pt, below=0.6cm of b2, xshift=1.4cm, text width=9.5cm, align=center] {
    \textcolor{red}{\textbf{Forensic Red Flag:}} Dishonest management conceals deteriorating customer solvency by issuing manual credit notes to cancel old unpaid invoices and reissuing fresh invoices (``aging resetting''), keeping balances artificially inside the $0$--$30$ day bucket!
};

\end{tikzpicture}
\caption{Accounts Receivable Aging Transition and Provisioning Escalation.}
\label{fig:ar_aging_ladder}
\end{figure}

% ----------------------------------------------------------------------
% PART 4: DEEP INTUITION & CONCEPTUAL DERIVATIONS
% ----------------------------------------------------------------------
\section{Deep Intuition \& Conceptual Derivations}

\subsection{The Masking Effect in Parametric Z-Score Audits}
Why does a standard $Z$-score fail in forensic fraud detection?
Suppose an enterprise ledger contains $1{,}000$ routine supplier payments averaging Rs.~$10{,}000$ with standard deviation $s = 2{,}000$. An embezzling controller introduces a massive single fraudulent payment of Rs.~$5{,}00{,}00{,}000$ (Rs.~$5$ Crores):
\begin{align*}
\bar{x}_{\text{new}} &\approx \frac{1000 \times 10{,}000 + 50{,}000{,}000}{1001} \approx \text{Rs.~} 59{,}940 \\
s_{\text{new}} &\approx \sqrt{\frac{1}{1000} \left(1000 \times (10{,}000 - 59{,}940)^2 + (50{,}000{,}000 - 59{,}940)^2\right)} \approx \text{Rs.~} 15{,}79{,}000
\end{align*}
Now, evaluate another suspicious payment of Rs.~$20{,}00{,}000$ (Rs.~$20$ Lakhs):
\begin{equation}
Z_i = \frac{2{,}000{,}000 - 59{,}940}{1{,}579{,}000} \approx 1.23
\end{equation}
A payment $200\times$ larger than normal operations yields a $Z$-score of only $1.23$, completely escaping detection because the massive Rs.~$5$ Crore transaction single-handedly inflated the sample standard deviation by $800\times$! This is known in robust statistics as the \textbf{masking effect}. 

Under the Modified $Z$-score:
\begin{align*}
\tilde{x} &= 10{,}000, \quad MAD = 1{,}349 \\
M_i &= \frac{0.6745 \times (2{,}000{,}000 - 10{,}000)}{1{,}349} \approx 995.2 \gg 3.5 \implies \text{\textcolor{red}{\textbf{Instantly Flagged!}}}
\end{align*}
The Median Absolute Deviation remains impervious to the extreme value (breakdown point of $50\%$).

% ----------------------------------------------------------------------
% PART 5: EXPLICIT GOTCHAS & DIAGNOSTIC MISCONCEPTIONS
% ----------------------------------------------------------------------
\section{Explicit Gotchas \& Diagnostic Misconceptions}

\begin{trapbox}{Gotcha 1: Conflating IRR with Actual Project Return}
\textbf{The Trap:} Assuming that an approved capital project with an $IRR$ of $35\%$ will compound the firm's invested capital at $35\%$ over its lifetime.\\
\textbf{Forensic Reality:} The mathematical formulation of $IRR$ inherently assumes that all intermediate project cash inflows can be reinvested at the $IRR$ itself ($35\%$), which is commercially unrealistic. The Modified Internal Rate of Return ($MIRR$), which reinvests inflows at the realistic Weighted Average Cost of Capital ($WACC$), provides the true economic yield.
\end{trapbox}

\begin{trapbox}{Gotcha 2: The Parametric Z-Score Breakdown in Skewed Financial Data}
\textbf{The Trap:} Applying Gaussian $\mu \pm 3\sigma$ filters to procurement payment ledgers and concluding that transactions below $3\sigma$ are normal.\\
\textbf{Forensic Reality:} Commercial transaction ledgers follow heavy-tailed distributions (log-normal or Pareto power-laws). In heavy-tailed data, extreme outliers distort both $\bar{x}$ and $\sigma$, causing severe masking where secondary fraudulent payments hide in plain sight.
\end{trapbox}

\begin{trapbox}{Gotcha 3: Neglecting the Low-Value Fraud Tail in Pareto Stratification}
\textbf{The Trap:} Auditing only the top $20\%$ of high-value disbursements under the Pareto 80/20 rule, assuming the remaining $80\%$ is immaterial.\\
\textbf{Forensic Reality:} Fraudsters are acutely aware of internal audit sampling thresholds. In procurement embezzlement, corrupt employees repeatedly submit thousands of low-value fake invoices (e.g., Rs.~$4{,}950$ to avoid a Rs.~$5{,}000$ voucher signoff rule). Auditing only top-tier transactions misses structuring and smurfing schemes entirely.
\end{trapbox}

\begin{trapbox}{Gotcha 4: Reversing Aging Provisions to Manufacture Paper Profits}
\textbf{The Trap:} Believing that changes in the bad debt provision matrix only reflect improved customer creditworthiness.\\
\textbf{Forensic Reality:} Dishonest management routinely reduces provision percentages on aged buckets (e.g., cutting the $>90$ days allowance from $50\%$ to $15\%$) during earnings slumps. This artificial adjustment reduces current-period bad debt expense, inflating reported Net Income without a single rupee of cash collection.
\end{trapbox}

% ----------------------------------------------------------------------
% PART 6: COMPARATIVE SYNTHESIS & SUMMARY TABLE
% ----------------------------------------------------------------------
\section{Comparative Synthesis \& Summary Table}

Table~\ref{tab:comparison_outliers} compares the statistical properties, breakdown bounds, and forensic suitability of parametric, non-parametric, and robust outlier detection algorithms.

\begin{table}[htbp]
\centering
\small
\begin{tabularx}{\textwidth}{lXXX}
\toprule
\textbf{Dimension} & \textbf{Parametric $Z$-Score} & \textbf{Tukey Boxplot (IQR)} & \textbf{Modified $Z$-Score (MAD)} \\
\midrule
\textbf{Distribution Assumption} & Strict Gaussian (Normal) distribution required. & Distribution-free; semi-parametric rank ordering. & Assumes underlying symmetric distribution around median. \\
\addlinespace
\textbf{Breakdown Point} & $0\%$ (a single extreme outlier can distort $\bar{x}$ and $s$). & $25\%$ (requires $>25\%$ contamination to corrupt quartiles). & $50\%$ (highest possible resistance to contamination). \\
\addlinespace
\textbf{Vulnerability to Masking} & Severe: high-value fraud inflates $s$, masking other anomalies. & Low: inner/outer fences remain stable under extreme spikes. & Minimal: median and MAD are completely robust to extreme tails. \\
\addlinespace
\textbf{Detection Threshold} & $|Z_i| > 3.0$ standard deviations. & $x_i > Q_3 + 1.5 \cdot IQR$ (outlier); $> Q_3 + 3.0 \cdot IQR$ (extreme). & $|M_i| > 3.5$ modified standard score. \\
\addlinespace
\textbf{Primary Audit Application} & Controlled manufacturing physical tolerances; standardized fees. & Accounts Receivable invoice aging; travel expense distributions. & Procurement journal audits; vendor payment amounts; wire transfers. \\
\bottomrule
\end{tabularx}
\caption{Comparative Matrix of Statistical Outlier Detection Frameworks in Financial Forensics.}
\label{tab:comparison_outliers}
\end{table}

% ----------------------------------------------------------------------
% PART 7: PYTHON / SCIENTIFIC COMPUTING LAB
% ----------------------------------------------------------------------
\section{Python / Scientific Computing Lab: Automated AR Aging \& Robust Outlier Detection}

Listing~\ref{lst:aging_outlier_engine} provides a complete, production-ready Python program that constructs accounts receivable aging schedules, calculates expected credit loss provisions, and executes robust Modified $Z$-score outlier filtering.

\begin{lstlisting}[language=Python, caption={Automated Invoice Aging \& Robust Outlier Detection in Python}, label={lst:aging_outlier_engine}]
import pandas as pd
import numpy as np

def audit_accounts_receivable(invoices_df: pd.DataFrame, as_of_date: str = '2025-03-31') -> dict:
    """
    Constructs an Accounts Receivable Aging Schedule and performs 
    Modified Z-score outlier detection on outstanding invoice amounts.
    
    Parameters:
        invoices_df (pd.DataFrame): Columns ['inv_id', 'customer', 'inv_date', 'amount']
        as_of_date (str): Audit balance cut-off date ('YYYY-MM-DD')
        
    Returns:
        dict: Aging schedule table, provisioning summary, and flagged outliers.
    """
    df = invoices_df.copy()
    df['inv_date'] = pd.to_datetime(df['inv_date'])
    ref_date = pd.to_datetime(as_of_date)
    
    # 1. Compute Days Past Due
    df['days_outstanding'] = (ref_date - df['inv_date']).dt.days
    
    # 2. Stratify into Aging Buckets
    bins = [-np.inf, 30, 60, 90, np.inf]
    labels = ['0-30 Days', '31-60 Days', '61-90 Days', '>90 Days']
    df['aging_bucket'] = pd.cut(df['days_outstanding'], bins=bins, labels=labels)
    
    # 3. Provisioning Rates Matrix
    prov_rates = {'0-30 Days': 0.01, '31-60 Days': 0.05, '61-90 Days': 0.20, '>90 Days': 0.50}
    df['provision_rate'] = df['aging_bucket'].map(prov_rates).astype(float)
    df['required_provision'] = df['amount'] * df['provision_rate']
    
    # Aging Summary Table
    aging_summary = df.groupby('aging_bucket', observed=False).agg(
        invoice_count=('inv_id', 'count'),
        total_balance=('amount', 'sum'),
        total_provision=('required_provision', 'sum')
    ).reset_index()
    aging_summary['effective_provision%'] = (
        aging_summary['total_provision'] / aging_summary['total_balance'] * 100.0
    ).round(2)
    
    # 4. Robust Outlier Detection: Modified Z-Score
    med = df['amount'].median()
    mad = np.median(np.abs(df['amount'] - med))
    if mad == 0:
        mad = np.mean(np.abs(df['amount'] - med))  # Fallback to mean absolute deviation
        
    df['modified_z'] = 0.6745 * (df['amount'] - med) / mad
    outliers = df[df['modified_z'].abs() > 3.5].copy()
    
    return {
        'total_ar_balance': df['amount'].sum(),
        'total_provision_required': df['required_provision'].sum(),
        'aging_schedule': aging_summary,
        'outlier_count': len(outliers),
        'flagged_outliers': outliers[['inv_id', 'customer', 'amount', 'days_outstanding', 'modified_z']]
    }

# -------------------------------------------------------------
# Test Harness with Aging Lag & Embezzlement Outliers
# -------------------------------------------------------------
if __name__ == '__main__':
    data = [
        {'inv_id': 'INV-101', 'customer': 'Alpha Corp', 'inv_date': '2025-03-15', 'amount': 15000.0},
        {'inv_id': 'INV-102', 'customer': 'Beta Ltd', 'inv_date': '2025-02-20', 'amount': 18500.0},
        {'inv_id': 'INV-103', 'customer': 'Gamma Ind', 'inv_date': '2025-01-10', 'amount': 22000.0},
        {'inv_id': 'INV-104', 'customer': 'Delta Retail', 'inv_date': '2024-11-05', 'amount': 35000.0},
        {'inv_id': 'INV-105', 'customer': 'Epsilon Tech', 'inv_date': '2025-03-25', 'amount': 14200.0},
        {'inv_id': 'INV-106', 'customer': 'Shell Co X', 'inv_date': '2025-03-28', 'amount': 850000.0} # Outlier!
    ]
    df_inv = pd.DataFrame(data)
    audit = audit_accounts_receivable(df_inv, as_of_date='2025-03-31')
    
    print(f"Total AR Outstanding: Rs. {audit['total_ar_balance']:,.2f}")
    print(f"Total Required Provision: Rs. {audit['total_provision_required']:,.2f}\n")
    print("Accounts Receivable Aging Schedule:")
    print(audit['aging_schedule'].to_string(index=False))
    print(f"\nFlagged Outliers (|Mod-Z| > 3.5): {audit['outlier_count']}")
    print(audit['flagged_outliers'].to_string(index=False))
\end{lstlisting}

% ----------------------------------------------------------------------
% PART 8: COMPREHENSIVE PRACTICE PROBLEMS
% ----------------------------------------------------------------------
\section{Comprehensive Practice Problems}
\textit{Complete mathematical derivations and diagnostic MCQ distractor rationales are located in the Appendix.}

\subsection*{Part A: Analytical \& Computational Problems}

\begin{enumerate}[label=\textbf{P6.\arabic*}]
    \item \textbf{Capital Budgeting Valuation \& Hurdle Rate Arbitrage:}\\
    An infrastructure developer is evaluating an industrial solar generation plant requiring an initial capital expenditure of Rs.~$100$ Crores at $t = 0$. Expected future net cash flows are:
    \begin{itemize}
        \item Year 1: Rs.~$25$ Crores, Year 2: Rs.~$35$ Crores, Year 3: Rs.~$40$ Crores, Year 4: Rs.~$45$ Crores.
    \end{itemize}
    \begin{enumerate}[label=(\alph*)]
        \item Compute the project's Net Present Value ($NPV$) assuming an enterprise cost of capital $k = 10\%$.
        \item State the decision rule for $NPV$ and determine whether the investment should be approved.
        \item If management artificially depresses the hurdle rate to $k = 6\%$ to push through a favored project, quantify the resulting artificial expansion in reported $NPV$.
    \end{enumerate}

    \item \textbf{Accounts Receivable Aging Schedule \& Provisioning Forensics:}\\
    At the close of the financial year, a logistics enterprise presents the following gross accounts receivable balances across aging brackets:
    \begin{center}
    \begin{tabular}{lrr}
    \toprule
    \textbf{Aging Bracket} & \textbf{Gross Receivables} & \textbf{Historical Default Rate} \\
    \midrule
    $0$--$30$ Days & Rs.~$50{,}00{,}000$ & $1.5\%$ \\
    $31$--$60$ Days & Rs.~$20{,}00{,}000$ & $6.0\%$ \\
    $61$--$90$ Days & Rs.~$12{,}00{,}000$ & $25.0\%$ \\
    $>90$ Days & Rs.~$8{,}00{,}000$ & $60.0\%$ \\
    \bottomrule
    \end{tabular}
    \end{center}
    \begin{enumerate}[label=(\alph*)]
        \item Calculate the total expected credit loss allowance required on the balance sheet.
        \item To meet an earnings benchmark, the CFO reclassifies Rs.~$6{,}00{,}000$ of receivables from the $>90$ Days bucket into the $0$--$30$ Days bucket by issuing replacement invoices. Compute the fraudulent reduction in bad debt expense and the percentage overstatement of pre-tax profit.
    \end{enumerate}

    \item \textbf{The Masking Effect in Forensic Outlier Detection:}\\
    A transaction ledger contains $n = 5$ vendor disbursements (in Rs.~Thousands): $[10, 12, 11, 13, 500]$.
    \begin{enumerate}[label=(\alph*)]
        \item Compute the sample mean $\bar{x}$ and sample standard deviation $s$.
        \item Compute the standard $Z$-score for the value $x = 13$.
        \item Compute the sample median $\tilde{x}$ and the Median Absolute Deviation ($MAD$).
        \item Compute the Modified $Z$-score $M_i$ for $x = 13$ and $x = 500$. Demonstrate why the Modified $Z$-score avoids the masking effect.
    \end{enumerate}

    \item \textbf{Tukey Interquartile Range (IQR) Non-Parametric Boundaries:}\\
    An auditor reviews the following ordered sample of monthly procurement office supply claims (in Rs.):
    \begin{equation*}
    [1200, 1400, 1500, 1650, 1800, 1950, 2100, 2300, 2800, 3100, 4500, 9800]
    \end{equation*}
    \begin{enumerate}[label=(\alph*)]
        \item Determine the first quartile ($Q_1$), median ($Q_2$), third quartile ($Q_3$), and $IQR$.
        \item Calculate the inner fences $[Q_1 - 1.5 \cdot IQR, Q_3 + 1.5 \cdot IQR]$ and outer fences $[Q_1 - 3.0 \cdot IQR, Q_3 + 3.0 \cdot IQR]$.
        \item Identify which claims qualify as mild outliers and which qualify as extreme outliers.
    \end{enumerate}

    \item \textbf{Pareto 80/20 Stratification \& Smurfing Audit Design:}\\
    A commercial bank processes $100{,}000$ international payments totaling Rs.~$500$ Crores annually. Explain how an investigative auditor should stratify this dataset using the Pareto principle to inspect both high-value concentration risks and structured smurfing schemes.
\end{enumerate}

\subsection*{Part B: Diagnostic Multiple-Choice Questions}

\begin{enumerate}[label=\textbf{MCQ 6.\arabic*}]
    \item Why is the Modified $Z$-score utilizing Median Absolute Deviation ($MAD$) preferred over the classical Gaussian $Z$-score when auditing commercial transaction ledgers for fraudulent disbursements?
    \begin{enumerate}[label=(\alph*)]
        \item The Gaussian $Z$-score can only be calculated for datasets with fewer than $30$ observations.
        \item Extreme fraudulent payments dramatically inflate both the sample mean and standard deviation, artificially reducing the calculated $Z$-scores of secondary fraudulent transactions (the masking effect).
        \item The Modified $Z$-score guarantees that the sum of all transaction amounts equals zero.
        \item Microsoft Excel is incapable of computing standard deviations on financial columns.
    \end{enumerate}

    \item Under the International Financial Reporting Standard (IFRS 9) expected credit loss framework, an enterprise discovers that a major customer's credit rating was downgraded from investment grade to sub-prime, and invoices have aged past $90$ days. What accounting entry is mandatory?
    \begin{enumerate}[label=(\alph*)]
        \item Derecognize the customer's accounts receivable and book a direct charge to Share Capital.
        \item Transfer the receivable balance into Intangible Goodwill.
        \item Recognize a lifetime expected credit loss allowance through a debit to Impairment Loss (Profit \& Loss) and a credit to Allowance for Doubtful Accounts (Contra-Asset).
        \item Defer recognizing any loss until the customer formally files for bankruptcy liquidation in court.
    \end{enumerate}
\end{enumerate}
